% Source: source/普通高考/2007/2007新课标理(海南,宁夏).pdf, page 1, question 5.
% pdfimages -list returned no embedded images; redraw from the rendered vector chart.
% Evidence: tmp/review_2007_new_curriculum_science/grid/page1_grid100.jpg,
%   tmp/review_2007_new_curriculum_science/grid/page1_grid50.jpg,
%   tmp/review_2007_new_curriculum_science/crops/q5_full.png (no-grid crop).
% Elements: rounded 开始/结束 terminals; process rectangles k=1, S=0,
% S=S+2k, k=k+1; decision k<=50?; output parallelogram 输出 S.
% Complete node -> directed-edge list:
% 开始 -> k=1 -> S=0 -> k<=50?;
% 是 -> S=S+2k -> k=k+1 -> left loop-back arrow entering the shared
% decision-entry stem;
% 否 -> 输出 S -> 结束. All borders/connectors are solid; no dashed edges.
% Every frame node is centered with local parindent=0pt and compact padding.

\begin{tikzpicture}[
  baseline,
  >=Stealth,
  x=1cm,
  y=1cm,
  line width=0.5pt,
  line cap=round,
  line join=round,
  font=\normalsize,
  flowtext/.style={anchor=center,align=center,inner sep=0pt,
    execute at begin node={\setlength{\parindent}{0pt}}},
  flowbox/.style={draw,flowtext,minimum height=0.72cm,inner xsep=4pt,inner ysep=2pt},
  terminal/.style={flowbox,rounded corners=.18cm,minimum width=1.35cm,
    minimum height=.74cm},
  process/.style={flowbox},
  io/.style={flowbox,shape=trapezium,trapezium left angle=72,
    trapezium right angle=108,minimum height=.80cm},
  decision/.style={flowbox,shape=diamond,aspect=2.9,minimum width=4.10cm,
    minimum height=1.05cm,inner sep=0pt},
  branchlabel/.style={flowtext}
]
  \node[terminal] (start) at (0,8.75) {开始};
  \node[process,minimum width=1.55cm] (kinit) at (0,7.50) {$k=1$};
  \node[process,minimum width=1.55cm] (sinit) at (0,6.18) {$S=0$};
  \node[decision] (test) at (0,4.65) {$k\leqslant 50?$};
  \node[process,minimum width=3.30cm] (sum) at (0,3.18) {$S=S+2k$};
  \node[process,minimum width=2.55cm] (inc) at (0,1.72) {$k=k+1$};
  \node[io,minimum width=2.25cm] (output) at (3.20,3.18) {输出 $S$};
  \node[terminal] (finish) at (3.20,1.72) {结束};

  \draw[->] (start.south) -- (kinit.north);
  \draw[->] (kinit.south) -- (sinit.north);
  \draw[->] (sinit.south) -- (test.north);
  \draw[->] (test.south) -- node[branchlabel,right=2pt] {是} (sum.north);
  \draw[->] (sum.south) -- (inc.north);
  \draw[->] (inc.south) -- (0,0.85) -- (-2.70,0.85) -- (-2.70,5.55) -- (0,5.55);
  \draw (0,5.55) -- (test.north);

  \draw (test.east) -- node[branchlabel,above=2pt] {否} (3.20,4.65)
    -- (3.20,3.72);
  \draw[->] (3.20,3.72) -- (output.north);
  \draw[->] (output.south) -- (finish.north);

  \path[use as bounding box] (-3.05,0.55) rectangle (4.45,9.10);
\end{tikzpicture}
